% Template for reports/latex/naca_lewis_16in_ramjet.py, filled by tools/latex_report.py.
% Every number below comes from the script, not from this file.
\documentclass{cadstudio}
\title{\VAR{title}}
\subtitle{\VAR{subtitle}}
\documentid{\VAR{document_id}}
\revision{\VAR{revision}}
\date{\VAR{date}}

\begin{document}
\maketitle

\begin{abstract}
  The duct coordinates tabulated in the CAD model of the engine of reference~\cite{E51C16} are
  reduced to a flow-area distribution, and one-dimensional gas dynamics is applied at the two
  free-jet Mach numbers of that report and over the travel of the exit plug. The conical shock of
  the \VAR{"%.0f"|format(cone)}\unit{\degree} spike is found to lie within
  \VAR{"%.1f"|format(lip_gap)}\unit{\degree} of the line from the spike tip to the cowl lip at both
  Mach numbers, the ideal total-pressure recovery at $M_\infty = \VAR{"%.2f"|format(design.mach)}$ is
  \VAR{"%.3f"|format(design.spike)}, and a choked exit sets the chamber Mach number between
  \VAR{"%.2f"|format(plug[0].mach)} and \VAR{"%.2f"|format(plug[-1].mach)}. The results are ideal
  estimates derived from the model geometry; they are not test data.
\end{abstract}

\begin{nomenclature}
  \item[$A$] flow area
  \item[$A^*$] flow area at sonic conditions
  \item[$M$] Mach number
  \item[$p_0$, $T_0$] total pressure and temperature
  \item[$r_s$, $r_b$] shell and inner-body radii
  \item[$V_r$] radial velocity in conical flow, divided by the maximum speed
  \item[$x$] station from the inlet lip
  \item[$\beta$] shock angle
  \item[$\gamma$] ratio of specific heats
  \item[$\eta$] total-pressure recovery
  \item[$\theta$] polar angle from the cone axis
  \item[$c$, $e$] subscripts: combustion chamber, nozzle exit
  \item[$n$] subscript: component normal to the shock
  \item[$1$, $2$] subscripts: upstream, downstream of a shock
  \item[$\infty$] subscript: free stream
\end{nomenclature}

\section{Introduction}

\BLOCK{if render}
The CAD model (figure~\ref{fig:cad}) reconstructs the 16-inch ramjet tested in a free jet at the
NACA Lewis laboratory~\cite{E51C16,E52D08}.
\BLOCK{else}
The CAD model reconstructs the 16-inch ramjet tested in a free jet at the NACA Lewis
laboratory~\cite{E51C16,E52D08}.
\BLOCK{endif}
Its shell, centre-body and spike profiles are the tabulated coordinates of
reference~\cite{E51C16}. This note checks those coordinates against elementary gas dynamics: the
area distribution of the duct, the position of the spike shock relative to the cowl lip, and the
chamber Mach number implied by the exit plug.

\BLOCK{if render}
\begin{figure}[!htb]
  \includegraphics[width=0.9\textwidth]{\VAR{render}}
  \caption{Exterior of the CAD assembly. Flow runs from the spike, lower right, to the nozzle.}
  \label{fig:cad}
\end{figure}
\BLOCK{endif}

\section{Flow path}

Stations $x$ are measured from the inlet lip along the flow axis. The shell and inner-body radii
$r_s(x)$ and $r_b(x)$ are linear interpolations of the model tables, with the spike at its
$M_\infty = \VAR{"%.2f"|format(design.mach)}$ position, and the annular flow area is
\begin{equation}
  A(x) = \pi\left[r_s(x)^2 - r_b(x)^2\right]. \label{eq:area}
\end{equation}
Figure~\ref{fig:duct} shows the result and table~\ref{tab:stations} lists it at the main stations.
The cowl lip encloses a capture area of \qty{\VAR{"%.0f"|format(capture_cm2)}}{\centi\metre\squared}
and the chamber area is \qty{\VAR{"%.0f"|format(chamber_cm2)}}{\centi\metre\squared}. The annulus
grows by a factor of \VAR{"%.2f"|format(diffuser_ratio)} between the lip and the station where the
shell reaches the chamber diameter. With the plug clear, the nozzle exit is
\VAR{"%.0f"|format(exit_fraction * 100)}\,\% of the chamber area, which coincides with the upper end
of the plug range given in reference~\cite{E51C16}.

\begin{figure}[!htb]
  \includegraphics{duct.pdf}
  \caption{Shell and inner-body radii from the model tables, and the flow area of
    equation~\eqref{eq:area}. The tail plug and the flame-holder blockage are not included.}
  \label{fig:duct}
\end{figure}

\begin{table}[!htb]
  \caption{Flow area at the main stations}
  \label{tab:stations}
  \begin{tabular}{@{}l S[table-format=4.1] S[table-format=3.1] S[table-format=3.1]
      S[table-format=4.0] S[table-format=1.3]@{}}
    \toprule
    Station & {$x$, mm} & {$r_s$, mm} & {$r_b$, mm} & {$A$, \unit{cm^2}} & {$A/A_c$} \\
    \midrule
\BLOCK{for row in stations}
    \VAR{row.label} & \VAR{"%.1f"|format(row.x_mm)} & \VAR{"%.1f"|format(row.shell_mm)}
      & \VAR{"%.1f"|format(row.body_mm)} & \VAR{"%.0f"|format(row.area_cm2)} & \VAR{"%.3f"|format(row.ratio)} \\
\BLOCK{endfor}
    \bottomrule
  \end{tabular}
  \tablenote{The centre body is hollow; pilot flow inside it is not counted.}
\end{table}

\section{Inlet shock system}

For a perfect gas with $\gamma = \VAR{gamma_air}$ the free-stream total conditions
are~\cite{AMES1135}
\begin{equation}
  \frac{T_0}{T_\infty} = 1 + \frac{\gamma - 1}{2} M_\infty^2, \qquad
  \frac{p_0}{p_\infty} = \left(\frac{T_0}{T_\infty}\right)^{\gamma/(\gamma - 1)}. \label{eq:total}
\end{equation}
A shock with upstream normal Mach number $M_n$ keeps the fraction
\begin{equation}
  \frac{p_{02}}{p_{01}} =
  \left[\frac{(\gamma + 1) M_n^2}{(\gamma - 1) M_n^2 + 2}\right]^{\gamma/(\gamma - 1)}
  \left[\frac{\gamma + 1}{2 \gamma M_n^2 - (\gamma - 1)}\right]^{1/(\gamma - 1)} \label{eq:shock}
\end{equation}
of the total pressure, with $M_n = M_\infty \sin\beta$ for a shock inclined at $\beta$ to the flow.
The spike is a cone of \VAR{"%.0f"|format(cone)}\unit{\degree} half-angle. Its shock angle follows
from the Taylor--Maccoll equation~\cite{TM1933} for the radial velocity $V_r(\theta)$ of the conical
flow, with $V_r' = \mathrm{d}V_r/\mathrm{d}\theta$:
\begin{equation}
  \frac{\gamma - 1}{2}\left[1 - V_r^2 - V_r'^2\right]\left[2 V_r + V_r' \cot\theta + V_r''\right]
  - V_r' \left[V_r V_r' + V_r' V_r''\right] = 0. \label{eq:tm}
\end{equation}
Equation~\eqref{eq:tm} is integrated from the shock to the ray where $V_r' = 0$, and $\beta$ is
iterated until that ray is the cone surface. The inlet recovery is then estimated as the conical
shock followed by a normal shock at the Mach number $M_2$ just behind it,
\begin{equation}
  \eta = \left.\frac{p_{02}}{p_{01}}\right|_{M_n = M_\infty \sin\beta} \times
         \left.\frac{p_{02}}{p_{01}}\right|_{M_n = M_2}. \label{eq:recovery}
\end{equation}

Table~\ref{tab:inlet} gives the results. The computed shock angles lie within
\VAR{"%.1f"|format(lip_gap)}\unit{\degree} of the line from the spike tip to the cowl lip at both
Mach numbers, so the two tip projections of reference~\cite{E51C16} place the conical shock close
to the lip in each case.

\begin{table}[!htb]
  \caption{Ideal inlet conditions at the two free-jet Mach numbers}
  \label{tab:inlet}
  \begin{tabular}{@{}l l \BLOCK{for c in conditions}S[table-format=3.3] \BLOCK{endfor}@{}}
    \toprule
    & & \multicolumn{\VAR{conditions|length}}{c}{$M_\infty$} \\
    \cmidrule(l){3-\VAR{conditions|length + 2}}
    Quantity & Basis
\BLOCK{for c in conditions}
      & {\VAR{"%.2f"|format(c.mach)}}
\BLOCK{endfor}
    \\
    \midrule
    Total temperature ratio, $T_0/T_\infty$ & Eq.~\eqref{eq:total}
      \BLOCK{for c in conditions}& \VAR{"%.3f"|format(c.t_ratio)} \BLOCK{endfor}\\
    Total pressure ratio, $p_0/p_\infty$ & Eq.~\eqref{eq:total}
      \BLOCK{for c in conditions}& \VAR{"%.3f"|format(c.p_ratio)} \BLOCK{endfor}\\
    Spike-tip projection, mm & Ref.~\cite{E51C16}
      \BLOCK{for c in conditions}& \VAR{"%.1f"|format(c.projection_mm)} \BLOCK{endfor}\\
    Tip-to-lip line angle, deg & Geometry
      \BLOCK{for c in conditions}& \VAR{"%.1f"|format(c.lip)} \BLOCK{endfor}\\
    Conical shock angle, $\beta$, deg & Eq.~\eqref{eq:tm}
      \BLOCK{for c in conditions}& \VAR{"%.1f"|format(c.shock)} \BLOCK{endfor}\\
    Mach number behind the shock, $M_2$ & Eq.~\eqref{eq:tm}
      \BLOCK{for c in conditions}& \VAR{"%.3f"|format(c.mach_2)} \BLOCK{endfor}\\
    Mach number on the cone surface & Eq.~\eqref{eq:tm}
      \BLOCK{for c in conditions}& \VAR{"%.3f"|format(c.surface_mach)} \BLOCK{endfor}\\
    Recovery, normal shock alone & Eq.~\eqref{eq:shock}
      \BLOCK{for c in conditions}& \VAR{"%.3f"|format(c.normal)} \BLOCK{endfor}\\
    Recovery, spike inlet, $\eta$ & Eq.~\eqref{eq:recovery}
      \BLOCK{for c in conditions}& \VAR{"%.3f"|format(c.spike)} \BLOCK{endfor}\\
    Recovery, MIL-E-5008B & Eq.~\eqref{eq:mil}
      \BLOCK{for c in conditions}& \VAR{"%.3f"|format(c.mil)} \BLOCK{endfor}\\
    \bottomrule
  \end{tabular}
\end{table}

Figure~\ref{fig:recovery} compares equation~\eqref{eq:recovery} with a single normal shock and with
the reference recovery of specification MIL-E-5008B,
\begin{equation}
  \eta_{\mathrm{MIL}} = 1 - 0.075 \left(M_\infty - 1\right)^{1.35}, \qquad 1 < M_\infty < 5. \label{eq:mil}
\end{equation}
The spike curve starts where the conical shock attaches. Viscous and subsonic-diffuser losses are
not included, so the curves are upper bounds for each shock arrangement.

\begin{figure}[!htb]
  \includegraphics{recovery.pdf}
  \caption{Ideal total-pressure recovery against free-stream Mach number. The symbols are the two
    free-jet conditions of table~\ref{tab:inlet}.}
  \label{fig:recovery}
\end{figure}

\section{Exit plug and chamber Mach number}

The movable tail plug sets the exit area $A_e$ between
\VAR{"%.0f"|format(plug_area[0] * 100)} and \VAR{"%.0f"|format(plug_area[1] * 100)}\,\% of the chamber
area~\cite{E51C16}. If the exit is choked and the flow from the chamber to the exit is isentropic,
the chamber Mach number $M_c$ is the subsonic root of the area--Mach relation~\cite{AMES1135} with
$A^* = A_e$:
\begin{equation}
  \frac{A_c}{A^*} = \frac{1}{M_c}
  \left[\frac{2}{\gamma + 1}\left(1 + \frac{\gamma - 1}{2} M_c^2\right)\right]^{(\gamma + 1)/[2(\gamma - 1)]}.
  \label{eq:areamach}
\end{equation}
Table~\ref{tab:plug} and figure~\ref{fig:plug} give $M_c$ at the end of the chamber, after heat
release, for an assumed $\gamma = \VAR{gamma_hot}$ in the combustion products and for cold air.
The two values of $\gamma$ change $M_c$ by at most \VAR{"%.3f"|format(gamma_gap)} over the plug
travel.

\begin{table}[!htb]
  \caption{Chamber Mach number for a choked exit}
  \label{tab:plug}
  \begin{tabular}{@{}S[table-format=1.2] S[table-format=1.3] S[table-format=1.3] S[table-format=1.3]@{}}
    \toprule
    & & \multicolumn{2}{c}{$M_c$} \\
    \cmidrule(l){3-4}
    {$A_e/A_c$} & {$A_c/A^*$} & {$\gamma = \VAR{gamma_hot}$} & {$\gamma = \VAR{gamma_air}$} \\
    \midrule
\BLOCK{for row in plug}
    \VAR{"%.2f"|format(row.fraction)} & \VAR{"%.3f"|format(row.ratio)} & \VAR{"%.3f"|format(row.mach)}
      & \VAR{"%.3f"|format(row.mach_air)} \\
\BLOCK{endfor}
    \bottomrule
  \end{tabular}
\end{table}

\begin{figure}[!htb]
  \includegraphics{plug.pdf}
  \caption{Chamber Mach number from equation~\eqref{eq:areamach} against exit-area fraction. The
    shaded band is the plug travel and the symbols are the rows of table~\ref{tab:plug}.}
  \label{fig:plug}
\end{figure}

\section{Assumptions}

\begin{enumerate}
  \item Air is a perfect gas with $\gamma = \VAR{gamma_air}$; $\gamma = \VAR{gamma_hot}$ for the
    combustion products is an assumed mean value.
  \item The flow is steady, inviscid and one-dimensional apart from the conical shock.
  \item Radii between tabulated stations are interpolated linearly, as in the CAD model.
  \item The flame holder, fuel bars, struts and tail plug are left out of the area distribution.
  \item The terminal normal shock is taken at $M_2$; the lower Mach number on the cone surface
    would give a slightly higher recovery.
\end{enumerate}

{\raggedright\noindent Geometry was read from \texttt{\VAR{source_path}}, SHA-256 digest
\texttt{\VAR{source_hash}}\dots\par}

\begin{thebibliography}{9}\raggedright
\BLOCK{for reference in references}
  \bibitem{\VAR{reference.id}} \VAR{reference.title}\BLOCK{if reference.locator} \VAR{reference.locator}\BLOCK{endif}

\BLOCK{endfor}
\end{thebibliography}

\end{document}
